\section{Local orthogonality forces complete continuity}
\label{sec:localization}

The graph criterion requires complete continuity of a free-space
linearization. The following theorem obtains it from a local
coordinate condition on the original Lipschitz map. The neighborhoods
are chosen around an arbitrary \emph{finite family of distinct
points}, and the orthogonality is required only after a common tail
cutoff. Neither a global partition nor uniform tail decay on bounded
sets is assumed.

\begin{theorem}[Localized linearization]
\label{thm:localized-linearization}
Let \((M,\rho,0)\) be a complete separable pointed metric space, and
let
\[
 U=\left(\bigoplus_{n\ge1}U_n\right)_2
\]
be a real Hilbert space with finite-dimensional blocks. Fix an
orthonormal basis in every \(U_n\), and write \(\Pi_{\ge l}\) for
projection onto \(U_{\ge l}=(\bigoplus_{n\ge l}U_n)_2\).
Suppose \(D:M\to U\) is Lipschitz with \(D(0)=0\), and that for
every \(m\ge1\) and every distinct \(z_1,\ldots,z_m\in M\) there are
open neighborhoods \(O_j\) of \(z_j\), an integer \(l\ge1\), and
mutually orthogonal coordinate subspaces \(H_j\subset U_{\ge l}\)
such that
\begin{equation}
\label{eq:local-orthogonality}
 \Pi_{\ge l}D(O_j)\subset H_j\qquad(1\le j\le m).
\end{equation}
A coordinate subspace means the closed span of a subset of the fixed
orthonormal basis. Then
\[
 q=\widehat D:\F(M)\to U
\]
is completely continuous.
\end{theorem}

The proof tests $D$ against weakly null unit vectors. After passage to
a subsequence, only finitely many points can retain local Lipschitz
constants above a fixed positive threshold. We first prove an
independent metric lemma that controls tests concentrated in shrinking
balls around those points. Compact reduction will then control the
remainder.

\begin{lemma}[Tests in shrinking balls]
\label[lemma]{lem:shrinking-tests}
Let \((M,\rho,0)\) be a pointed metric space, let \(z\in M\), and
let \(\mu_i\to0\) weakly in \(\F(M)\). Suppose \(r_i>0\),
\(r_i\to0\), and \(f_i:M\to\R\) satisfy
\[
 f_i(z)=0,\qquad \Lip(f_i)\le C,\qquad
 f_i(x)=0\quad\hbox{when }\rho(x,z)\ge r_i,
\]
where \(C\) does not depend on \(i\). Then
\(\langle f_i,\mu_i\rangle\to0\).
\end{lemma}

\begin{proof}
For \(s>0\), define
\[
 \phi_{z,s}(x)=\max\{0,\min\{1,\,2-2\rho(x,z)/s\}\}.
\]
This cutoff is one on \(B(z,s/2)\), zero outside \(B(z,s)\), takes
values in \([0,1]\), and has Lipschitz constant at most \(2/s\).
Whenever \(f(z)=0\) and \(\Lip(f)\le C\), we claim
\begin{equation}
\label{eq:cutoff-products}
 \Lip(f\phi_{z,s})\le3C,\qquad
 \Lip(f(1-\phi_{z,s}))\le3C.
\end{equation}
Let \(\psi\) be either cutoff in this claim. If
\(\psi(x)=\psi(y)\), the difference of the products is at most
\(C\rho(x,y)\). Otherwise at least one endpoint lies in \(B(z,s)\);
call it \(x\). Then \(|f(x)|\le Cs\), and
\[
 f(x)\psi(x)-f(y)\psi(y)
 =\psi(y)(f(x)-f(y))+f(x)(\psi(x)-\psi(y))
\]
has absolute value at most
\(C\rho(x,y)+Cs(2/s)\rho(x,y)=3C\rho(x,y)\).
Thus \eqref{eq:cutoff-products} needs no global bound on \(f\).

Suppose the conclusion fails. Passing to a subsequence, assume
\[
 |\langle f_i,\mu_i\rangle|\ge\epsilon>0\qquad(i\ge1).
\]
Then \(C>0\). We recursively select increasing indices \(i_j\),
numbers \(0<s_j<r_{i_j}\), and functions
\[
 h_j=f_{i_j}(1-\phi_{z,s_j})
\]
such that
\begin{equation}
\label{eq:annular-selection}
 |\langle h_j,\mu_{i_j}\rangle|>\epsilon/2,\qquad
 \Lip(h_j)\le3C,\qquad r_{i_{j+1}}<s_j/2.
\end{equation}
After choosing \(i_j\), approximate \(\mu_{i_j}\) by a finite sum of
evaluations
\[
 \nu_j=\sum_{k=1}^{N_j}a_{jk}\delta(x_{jk}),\qquad
 \|\mu_{i_j}-\nu_j\|<\frac{\epsilon}{6C}.
\]
Choose \(s_j\in(0,r_{i_j})\) smaller than every positive distance
from \(z\) to a member of
\(\{0,x_{j1},\ldots,x_{jN_j}\}\). If that finite set has no positive
such distances, impose only \(s_j<r_{i_j}\).
The function \(p_j=f_{i_j}\phi_{z,s_j}\) vanishes at every one of
these points: outside the ball the cutoff is zero, and at a point
equal to \(z\) the factor \(f_{i_j}\) is zero. Including the
basepoint is essential when \(z\ne0\). It gives
\[
 \langle p_j,\nu_j\rangle
 =\sum_{k=1}^{N_j}a_{jk}\bigl(p_j(x_{jk})-p_j(0)\bigr)=0.
\]
By \eqref{eq:cutoff-products},
\[
 |\langle p_j,\mu_{i_j}\rangle|
 \le3C\|\mu_{i_j}-\nu_j\|<\epsilon/2.
\]
Subtracting from the original pairing proves the first inequality in
\eqref{eq:annular-selection}; the second is
\eqref{eq:cutoff-products}. Finally choose \(i_{j+1}>i_j\) so large
that \(r_{i_{j+1}}<s_j/2\). This completes the recursion.

The nonzero set of \(h_j\) lies in
\[
 \{x:s_j/2<\rho(x,z)<r_{i_j}\}.
\]
These annuli are pairwise disjoint. Every finite signed sum
\(H=\sum_{j\in A}\sigma_jh_j\), with \(\sigma_j\in\{-1,1\}\),
therefore satisfies \(\Lip(H)\le6C\). Indeed, at a point at most
one summand is nonzero. If two points lie in the same nonzero set,
the \(3C\) bound for that summand applies. If they lie in different
nonzero sets, each of the two contributing functions vanishes at the
other point, and their \(3C\) bounds add. The cases where one or both
points lie in no nonzero set are included. Subtracting \(H(0)\)
does not change this Lipschitz bound.

For every \(\mu\in\F(M)\) and finite \(A\subset\N\), real signs give
\[
 \sum_{j\in A}|\langle h_j,\mu\rangle|
 =\max_{\sigma_j\in\{-1,1\}}
   \left|\left\langle\sum_{j\in A}\sigma_jh_j,\mu\right\rangle\right|
 \le6C\|\mu\|.
\]
Taking increasing finite sets shows that
\[
 T:\F(M)\to\ell_1,\qquad T\mu=(\langle h_j,\mu\rangle)_{j\ge1}
\]
is a bounded linear operator of norm at most \(6C\).
Thus \(T\mu_{i_j}\to0\) weakly in \(\ell_1\). By the Schur property
of \(\ell_1\), its norms tend to zero.
But its \(j\)th coordinate has absolute value greater than
\(\epsilon/2\) by \eqref{eq:annular-selection}, a contradiction.
\end{proof}

\begin{proof}[Proof of \Cref{thm:localized-linearization}]
Let \(E=\F(M)\) and \(L=\Lip(D)=\|q\|\). If \(L=0\), then
\(D=0\) because \(D(0)=0\), and \(q=0\). Assume \(L>0\).
If \(q\) were not completely continuous, after taking a subsequence
there would be \(\mu_i\to0\) weakly in \(E\) and
\(\|q\mu_i\|\ge\epsilon_0>0\) for all \(i\). Set
\[
 b_i=\frac{q\mu_i}{\|q\mu_i\|},\qquad
 d_i(x)=\langle D(x),b_i\rangle,\qquad
 M_0=\sup_i\|\mu_i\|<\infty.
\]
The unit vectors \(b_i\) are weakly null. For each \(b\in U\),
\[
 |\langle b_i,b\rangle|
 \le\epsilon_0^{-1}|\langle q\mu_i,b\rangle|\longrightarrow0,
\]
since a bounded linear map preserves weak convergence. Consequently
\begin{equation}
\label{eq:scalar-tests}
 d_i(0)=0,\qquad \Lip(d_i)\le L,\qquad
 \langle d_i,\mu_i\rangle=\|q\mu_i\|\ge\epsilon_0.
\end{equation}
The functions \(d_i\) tend uniformly to zero on every compact subset
of \(M\). They tend pointwise to zero by weak nullity of \(b_i\).
Given a compact set and a positive error, take a finite sufficiently
fine metric net; pointwise convergence on that net and the common
Lipschitz bound control all points of the compact set.

\medskip
\noindent\emph{Only finitely many points can retain a large local
Lipschitz constant.}
Choose a countable base \(\mathcal B\) for \(M\), which exists by
separability. Take the Lipschitz constant on a set of at most one
point to be zero. For each \(O\in\mathcal B\),
\(\Lip(d_i|_O)\in[0,L]\). Successively take subsequences for the
countably many \(O\), then take a diagonal subsequence. Relabeling,
we may assume that all limits
\[
 a_O=\lim_{i\to\infty}\Lip(d_i|_O)
 \qquad(O\in\mathcal B)
\]
exist. This subsequence is fixed before the next parameter is chosen.

Fix \(\eta>0\) and define
\[
 Z=\{z\in M:a_O\ge\eta\ \hbox{for every }O\in\mathcal B
                         \hbox{ containing }z\}.
\]
Take any finite family of distinct \(z_1,\ldots,z_m\in Z\).
Apply \eqref{eq:local-orthogonality} and shrink the resulting
neighborhoods to base neighborhoods \(O_j\) containing \(z_j\).
Let \(P_j\) be projection onto \(H_j\), and put
\(\Pi_{<l}=I-\Pi_{\ge l}\). For \(x,y\in O_j\), the tail
\(\Pi_{\ge l}(D(x)-D(y))\) lies in \(H_j\). Hence
\[
 |d_i(x)-d_i(y)|
 \le L\rho(x,y)\bigl(\|P_jb_i\|+\|\Pi_{<l}b_i\|\bigr).
\]
The first \(l-1\) blocks are finite dimensional. Weak nullity of
\(b_i\) makes each of its finitely many coordinates there tend to
zero, so \(\|\Pi_{<l}b_i\|\to0\). Since \(a_{O_j}\ge\eta\),
the preceding inequality gives
\[
 \liminf_{i\to\infty}\|P_jb_i\|\ge\eta/L.
\]
The \(P_j\) have mutually orthogonal ranges, and \(\|b_i\|=1\).
It follows that
\[
 m\frac{\eta^2}{L^2}
 \le \liminf_{i\to\infty}\sum_{j=1}^m\|P_jb_i\|^2
 \le1.
\]
The same conclusion holds for every finite family of distinct points
of \(Z\), so \(Z\) is finite and
\begin{equation}
\label{eq:exceptional-cardinality}
 |Z|\le L^2/\eta^2.
\end{equation}

\medskip
\noindent\emph{Remove small neighborhoods of those finitely many
points.}
Choose \(r_0>0\) so that the balls \(B(z,r)\), \(z\in Z\), are
pairwise disjoint for every \(0<r<r_0\). If \(Z\) has at most one
point, any \(r_0>0\) suffices. With the cutoffs used in
\Cref{lem:shrinking-tests}, put
\[
 f_{i,z,r}=\phi_{z,r}(d_i-d_i(z)),\qquad
 g_{i,r}=\sum_{z\in Z}f_{i,z,r},\qquad
 R_{i,r}=d_i-g_{i,r}.
\]
An empty sum is zero. Equation~\eqref{eq:cutoff-products} gives
\(\Lip(f_{i,z,r})\le3L\). Their nonzero sets lie in disjoint balls,
so the two-point argument for finite signed sums in the proof of
\Cref{lem:shrinking-tests} gives
\begin{equation}
\label{eq:global-remainder}
 \Lip(g_{i,r})\le6L,\qquad \Lip(R_{i,r})\le7L.
\end{equation}
These constants do not depend on \(r\) or on \(|Z|\).

For each fixed \(z\in Z\), we have the ordered limit
\begin{equation}
\label{eq:one-shrinking-limit}
 \lim_{r\downarrow0}\limsup_{i\to\infty}
      |\langle f_{i,z,r},\mu_i\rangle|=0.
\end{equation}
If this failed, some \(a>0\) and radii \(r_j\downarrow0\) would
satisfy
\(\limsup_i|\langle f_{i,z,r_j},\mu_i\rangle|\ge2a\).
Choose increasing \(i_j\) with
\(|\langle f_{i_j,z,r_j},\mu_{i_j}\rangle|\ge a\).
These tests vanish at \(z\), have Lipschitz constants at most \(3L\),
and vanish outside \(B(z,r_j)\). The weakly null subsequence
\(\mu_{i_j}\) and \Cref{lem:shrinking-tests} give a contradiction.
Because \(Z\) is finite, the triangle inequality and
\eqref{eq:one-shrinking-limit} imply
\begin{equation}
\label{eq:shrinking-limit}
 \lim_{r\downarrow0}\limsup_{i\to\infty}
      |\langle g_{i,r},\mu_i\rangle|=0.
\end{equation}
If \(Z\) is empty, \eqref{eq:shrinking-limit} is immediate.

The removed part is now controlled. It remains to show that the
remainder has small pairing. This is where completeness of \(M\)
enters through compact reduction.

\medskip
\noindent\emph{Estimate the remainder on one compact set.}
Fix \(\tau>0\). By \Cref{thm:compact-reduction}, choose a compact
\(K\subset M\), \(0\in K\), and \(\nu_i\in\F(K)\) with
\[
 \|\mu_i-\nu_i\|<\tau,\qquad
 \|\nu_i\|\le M_0+\tau.
\]
The set \(K\) depends on \(\tau\) and the sequence, but not on \(r\).
Now fix \(0<r<r_0\) while taking the sequence limit. Write \(m=|Z|\)
and
\[
 \Phi_r=\sum_{z\in Z}\phi_{z,r},\qquad
 \epsilon_i=\sup_{x\in K\cup Z}|d_i(x)|.
\]
The compact set \(K\cup Z\) contains \(0\), so
\(\epsilon_i\to0\), including when \(Z\) is empty.
Disjointness of the balls gives \(0\le\Phi_r\le1\), and
\begin{equation}
\label{eq:remainder-identity}
 R_{i,r}=(1-\Phi_r)d_i+\sum_{z\in Z}\phi_{z,r}d_i(z).
\end{equation}
For \(x\in K\) the right side is a convex combination of values
bounded in absolute value by \(\epsilon_i\). Therefore
\begin{equation}
\label{eq:remainder-sup}
 \sup_{x\in K}|R_{i,r}(x)|\le\epsilon_i.
\end{equation}
On each inner ball \(B(z,r/2)\), the remainder is the constant
\(d_i(z)\).

At any point \(p\notin Z\), choose \(O\in\mathcal B\) containing
\(p\) with \(a_O<\eta\). For \(x,y\in K\cap O\), expansion of
\eqref{eq:remainder-identity} gives
\begin{align*}
 R_{i,r}(x)-R_{i,r}(y)
 ={}&(1-\Phi_r(x))(d_i(x)-d_i(y))\\
 &+\sum_{z\in Z}
   \bigl(\phi_{z,r}(x)-\phi_{z,r}(y)\bigr)
   \bigl(d_i(z)-d_i(y)\bigr).
\end{align*}
The last factors have absolute value at most \(2\epsilon_i\), and
\(\Lip(\phi_{z,r})\le2/r\). Thus
\begin{equation}
\label{eq:local-remainder}
 \Lip(R_{i,r}|_{K\cap O})
 \le \Lip(d_i|_O)+\frac{4m}{r}\epsilon_i.
\end{equation}
The radius is fixed here, so the second term tends to zero with \(i\).

The inner balls at points of \(Z\), together with the chosen base
neighborhoods at points outside \(Z\), cover \(K\). Choose a finite
subcover \(V_1,\ldots,V_N\). The Lebesgue-number property supplies
\(\lambda>0\) such that any pair of points of \(K\) at distance less
than \(\lambda\) belongs to some one \(V_t\).
For an inner ball put \(c_{i,t}=0\); for a chosen base neighborhood
put
\[
 c_{i,t}=\Lip(d_i|_{V_t})+(4m/r)\epsilon_i.
\]
For pairs at distance less than \(\lambda\), the Lipschitz quotient
of \(R_{i,r}\) is at most \(\max_t c_{i,t}\), by constancy on inner
balls and \eqref{eq:local-remainder}. For pairs at distance at least
\(\lambda\), \eqref{eq:remainder-sup} bounds it by
\(2\epsilon_i/\lambda\). Therefore
\[
 \Lip(R_{i,r}|_K)
 \le\max\left\{\max_{1\le t\le N}c_{i,t},
                       \frac{2\epsilon_i}{\lambda}\right\}.
\]
If \(K\) has only one point, its Lipschitz constant is zero and the
same conclusion holds. There are finitely many \(V_t\), and for
each base neighborhood \(a_{V_t}<\eta\). Taking the upper limit gives
\begin{equation}
\label{eq:compact-remainder}
 \limsup_{i\to\infty}\Lip(R_{i,r}|_K)\le\eta.
\end{equation}

Use \eqref{eq:compact-pairing} for \(\nu_i\in\F(K)\), and use the
global bound \eqref{eq:global-remainder} for the approximation error.
We obtain
\begin{align*}
 |\langle R_{i,r},\mu_i\rangle|
 &\le |\langle R_{i,r},\nu_i\rangle|
       +|\langle R_{i,r},\mu_i-\nu_i\rangle|\\
 &\le (M_0+\tau)\Lip(R_{i,r}|_K)+7L\tau.
\end{align*}
With \(\eta,\tau,K,r\) fixed, \eqref{eq:compact-remainder} implies
\[
 \limsup_{i\to\infty}|\langle R_{i,r},\mu_i\rangle|
 \le \eta(M_0+\tau)+7L\tau.
\]
The right side is independent of \(r\). Since
\(d_i=R_{i,r}+g_{i,r}\), we may now let \(r\downarrow0\) and use
\eqref{eq:shrinking-limit} to get
\begin{equation}
\label{eq:final-local-bound}
 \limsup_{i\to\infty}|\langle d_i,\mu_i\rangle|
 \le \eta(M_0+\tau)+7L\tau
 \qquad\hbox{for every }\eta>0,\ \tau>0.
\end{equation}
The quantifier order is as follows: the diagonal subsequence is fixed
first; for each \(\eta\) choose \(Z,r_0\); for each \(\tau\) choose
\(K\); for a fixed \(r<r_0\) take \(i\to\infty\); then take
\(r\downarrow0\), \(\tau\downarrow0\), and finally
\(\eta\downarrow0\). No uniform convergence in both \(i\) and \(r\)
has been used. The last two limits in \eqref{eq:final-local-bound}
make its right side zero, contradicting \eqref{eq:scalar-tests}.
\end{proof}
